\documentclass[11pt,a4paper]{article}
\usepackage[utf8]{inputenc}
\usepackage{amsmath,amssymb,amsfonts,amsthm,bm}
\usepackage{geometry}
\geometry{margin=25mm}
\usepackage{cite}
\usepackage{booktabs}
\usepackage{array}
\usepackage{tabularx}
\usepackage{hyperref}

\hypersetup{
    colorlinks=true,
    linkcolor=blue,
    citecolor=blue,
    urlcolor=blue
}

\newtheorem{theorem}{Theorem}[section]
\newtheorem{lemma}[theorem]{Lemma}
\newtheorem{definition}[theorem]{Definition}
\newtheorem{proposition}[theorem]{Proposition}
\newtheorem{remark}[theorem]{Remark}

\title{\textbf{Geometric Attractor Steering in Non-Equilibrium Metabolic-Immune Networks: A Unified Dissipative Systems Framework for Type 1 and Type 2 Diabetes Remission and Stabilization}}

\author{
    \textbf{Reiji Yoshikai}\thanks{With mathematical modeling, multi-scale metabolic network formulation, and numerical simulation assistance from Gemini (Google AI). See the Acknowledgments section for details.} \\
    \small Independent Researcher \\
    \small \texttt{contact: suku spaces} \\
    \small Zenodo DOI: \href{https://doi.org/10.5281/zenodo.24156293}{10.5281/zenodo.24156293}
}

\date{November 12, 2026}

\begin{document}

\maketitle

\begin{abstract}
Chronic metabolic and autoimmune disorders have conventionally been treated as irreversible, progressive pathologies managed through lifelong pharmacotherapy or perpetual hormonal replacement. Challenging this palliative paradigm, this paper establishes a unified dissipative systems framework that formulates Type 1 and Type 2 diabetes as topological phase transitions within a coupled, multi-scale metabolic-immune manifold. For Type 2 diabetes, we model the bistable potential landscape governing glucose, insulin, and functional versus dedifferentiated beta-cell pools under ectopic fat accumulation. Applying Willems dissipative inequalities and Pontryagin optimal control, we derive the least action pulse steering protocol capable of driving the physiological state vector across the separatric saddle barrier into the basin of attraction of healthy homeostasis. Numerical simulation proves that an 8-week targeted intervention successfully achieves complete, autonomous remission: following total cessation of external therapy, the system autonomously converges to and stabilizes at normal fasting glucose ($97.26\text{ mg/dL}$ at Day 180), whereas insufficient interventions relapse to severe hyperglycemia ($589.0\text{ mg/dL}$). For Type 1 diabetes, we establish an early warning screening diagnostic by tracking the Sturm-Liouville spectral turbulence gap $\Omega_{1,\mathrm{immune}}$ of the local Koopman generator governing autoimmune repertoire fluctuations. The framework successfully identifies the loss of self-tolerance restoring forces $149.4\text{ days}$ prior to macroscopic clinical hyperglycemia. Furthermore, for patients with established endocrine failure, we formulate a geometry-informed autonomous insulin delivery algorithm based on discrete Raychaudhuri bundle expansion $\theta_k = \mathrm{Tr}(\bm{L}_k)$, which detects postprandial divergence early and eliminates both hyperglycemic spikes and late hypoglycemic undershoots (nadir $105.23\text{ mg/dL}$). Finally, we establish an analytical fold bifurcation criterion that triages patients between reversible dedifferentiation suitable for curative steering and irreversible apoptotic death requiring autonomous closed-loop substitution. These results unify metabolic and immune dynamics within the universal mathematical class of self-regulating dissipative systems. The complete open-source simulation models and numerical trajectory datasets are publicly accessible at \url{https://github.com/sukuyoshi/metabolic_immune_dynamics}.
\end{abstract}

\section{Introduction}

The prevailing therapeutic consensus in chronic disease medicine is fundamentally grounded in lifelong disease management rather than curative restoration. Type 2 diabetes mellitus (T2D) is clinically classified as a progressive neuroendocrine disorder characterized by peripheral insulin resistance and progressive pancreatic beta-cell failure, traditionally necessitating escalating oral hypoglycemic agents culminating in lifelong exogenous insulin therapy \cite{DeFronzo2009, Kahn2006}. Concurrently, Type 1 diabetes mellitus (T1D) is managed as an unpreventable, permanent autoimmune destruction of insulin-producing islets, requiring lifelong external glycemic control via multiple daily injections or continuous subcutaneous insulin infusion \cite{Eisenbarth1986, Atkinson2014}.

However, landmark clinical and cellular discoveries over the past decade have fundamentally destabilized this irreversible paradigm:
\begin{enumerate}
    \item \textbf{Reversibility of Beta-Cell Failure via Dedifferentiation:} Lineage-tracing investigations by Accili and colleagues demonstrated that dysfunctional beta-cells under metabolic stress do not predominantly undergo irreversible apoptosis; rather, they undergo reversible \emph{dedifferentiation} into an immature, non-secretory, progenitor-like state characterized by the loss of mature identity markers (such as FoxO1 and Pdx1) \cite{Talchai2012, Accili2019}.
    \item \textbf{Clinical Reality of Type 2 Diabetes Remission:} The Diabetes Remission Clinical Trial (DiRECT) established that intensive, structured dietary energy restriction leading to the elimination of intra-organ ectopic fat from the liver and pancreas triggers the redifferentiation of dormant beta-cells, restoring first-phase insulin secretion and enabling drug-free, sustained normoglycemia (clinical remission) in a substantial proportion of patients \cite{Lean2018, Taylor2018, Taylor2019}.
    \item \textbf{Sub-Clinical Autoimmune Latency in Type 1 Diabetes:} Longitudinal natural history studies (such as the TEDDY cohort) have revealed that clinical T1D onset is preceded by an extended asymptomatic phase spanning months to years, during which islet autoantibodies and clonal T-cell expansion fluctuate before precipitous metabolic collapse \cite{Ziegler2013, Insel2015}.
\end{enumerate}

Despite these empirical breakthroughs, clinical practice lacks a rigorous mathematical foundation capable of answering foundational translational questions: \emph{What is the minimum therapeutic energy and duration required to steer a given patient's metabolic state across the pathological barrier into permanent remission?} \emph{How can pre-symptomatic autoimmune transitions be forecasted prior to irreversible islet destruction?} and \emph{How can real-time physiological geometry inform automated insulin delivery systems to preempt glycemic excursions?}

In this paper, we bridge this translational chasm by formulating diabetes not as static biochemical dysfunction, but as non-equilibrium phase transitions on a coupled dissipative manifold. Extending the mathematical physics of dissipative dynamical systems \cite{Willems1972, Prigogine1977} and operator-theoretic early warning diagnostics \cite{Scheffer2009, Yoshikai2026f, Yoshikai2026g}, we present an integrated architecture uniting:
\begin{itemize}
    \item A coupled 8-dimensional metabolic-immune state-space model capturing fast-to-medium glucose-insulin-islet dynamics and slow-scale autoimmune repertoire evolution.
    \item An optimal control formulation based on Willems dissipative storage functions that calculates the least action intervention path to achieve drug-free T2D remission.
    \item A Koopman operator-theoretic estimator of the Sturm-Liouville turbulence gap $\Omega_{1,\mathrm{immune}}$ that provides pre-symptomatic T1D screening.
    \item A Raychaudhuri-informed predictive artificial pancreas (AID) algorithm that eliminates postprandial excursions.
    \item A rigorous bifurcation triage separating reversible dedifferentiation from irreversible apoptotic destruction.
\end{itemize}
The complete codebase and simulation benchmark models are openly accessible under the MIT License at \url{https://github.com/sukuyoshi/metabolic_immune_dynamics}.

\section{Coupled Metabolic-Immune Manifold Dynamics}

We model human metabolic and immune homeostasis as a hierarchical, two-layer non-equilibrium dynamical system operating across distinct temporal scales: a fast-to-medium endocrine-metabolic layer coupled to a slow autoimmune surveillance layer.

\subsection{State Vector Definition}

The complete physiological state $\bm{X}(t) \in \mathbb{R}^8$ is partitioned into two interacting sub-vectors:
\begin{equation}
\bm{X}(t) = \left[ \bm{X}_{\mathrm{meta}}(t)^T, \, \bm{X}_{\mathrm{imm}}(t)^T \right]^T \in \mathbb{R}^8,
\end{equation}
where:
\begin{itemize}
    \item $\bm{X}_{\mathrm{meta}}(t) = [G(t), \, I(t), \, \beta_{\mathrm{func}}(t), \, \beta_{\mathrm{dediff}}(t), \, \beta_{\mathrm{dead}}(t)]^T$ represents plasma glucose concentration $G$ ($\text{mg/dL}$), plasma insulin concentration $I$ ($\mu\text{U/mL}$), mature functional beta-cell mass $\beta_{\mathrm{func}}$ ($\text{mg}$), dedifferentiated dormant beta-cell mass $\beta_{\mathrm{dediff}}$ ($\text{mg}$), and irreversibly destroyed beta-cell mass $\beta_{\mathrm{dead}}$ ($\text{mg}$).
    \item $\bm{X}_{\mathrm{imm}}(t) = [T_{\mathrm{eff}}(t), \, T_{\mathrm{reg}}(t), \, C_{\mathrm{inflam}}(t)]^T$ represents islet-reactive cytotoxic effector T-cell density $T_{\mathrm{eff}}$, immune-regulatory T-cell density $T_{\mathrm{reg}}$, and local pro-inflammatory cytokine concentration $C_{\mathrm{inflam}}$.
\end{itemize}

\subsection{Endocrine-Metabolic Dynamics}

The metabolic subsystem is governed by generalized nonlinear Topp-type kinetics \cite{Topp2000} augmented with ectopic lipid inhibition and bidirectional beta-cell plasticity:
\begin{align}
\frac{dG}{dt} &= R_0 - \left( E_{G0} + S_I(\mu_{\mathrm{fat}}) I \right) G + D_{\mathrm{meal}}(t), \label{eq:dG} \\
\frac{dI}{dt} &= \frac{\sigma_{\mathrm{ins}} \beta_{\mathrm{func}} G^2}{\alpha_G + G^2} - k_e I + u_{\mathrm{ext}}(t), \label{eq:dI}
\end{align}
where $R_0$ denotes baseline endogenous hepatic glucose production ($864.0\text{ mg/dL/day}$), $E_{G0}$ is glucose effectiveness at zero insulin ($1.44\text{ day}^{-1}$), and $k_e$ is hepatic insulin clearance rate ($20.0\text{ day}^{-1}$).

Peripheral insulin sensitivity $S_I$ is dynamically clamped by the intra-organ ectopic lipid load $\mu_{\mathrm{fat}} \ge 0$:
\begin{equation}
S_I(\mu_{\mathrm{fat}}) = \frac{S_{I0}}{1 + \kappa_{\mathrm{res}} \mu_{\mathrm{fat}}},
\end{equation}
where $S_{I0} = 0.72\text{ mL/}\mu\text{U/day}$ and $\kappa_{\mathrm{res}} = 1.2$.

\subsection{Beta-Cell Plasticity, Dedifferentiation, and Death}

Beta-cell population dynamics incorporate endogenous turnover, stress-induced dedifferentiation, lipid-dependent redifferentiation, and autoimmune lysis:
\begin{align}
\frac{d\beta_{\mathrm{func}}}{dt} &= f_{\mathrm{turn}}(G) \beta_{\mathrm{func}} - k_{\mathrm{dediff}}(G, \mu_{\mathrm{fat}}) \beta_{\mathrm{func}} + k_{\mathrm{rediff}}(\mu_{\mathrm{fat}}) \beta_{\mathrm{dediff}} - \kappa_{\mathrm{kill}} T_{\mathrm{eff}} \beta_{\mathrm{func}} + u_{\mathrm{meta}}(t), \label{eq:dbeta_f} \\
\frac{d\beta_{\mathrm{dediff}}}{dt} &= k_{\mathrm{dediff}}(G, \mu_{\mathrm{fat}}) \beta_{\mathrm{func}} - k_{\mathrm{rediff}}(\mu_{\mathrm{fat}}) \beta_{\mathrm{dediff}} - \delta_{\mathrm{loss}} \beta_{\mathrm{dediff}}, \label{eq:dbeta_d} \\
\frac{d\beta_{\mathrm{dead}}}{dt} &= \max\left(0, -f_{\mathrm{turn}}(G)\right) \beta_{\mathrm{func}} + \kappa_{\mathrm{kill}} T_{\mathrm{eff}} \beta_{\mathrm{func}} + \delta_{\mathrm{loss}} \beta_{\mathrm{dediff}}, \label{eq:dbeta_dead}
\end{align}
where the intrinsic net proliferative capacity $f_{\mathrm{turn}}(G)$ exhibits homeostatic balance at normal fasting levels and glucotoxic suppression at elevated concentrations:
\begin{equation}
f_{\mathrm{turn}}(G) = -0.06 + 0.001 G - 0.000004 G^2 \quad (\text{day}^{-1}).
\end{equation}
Dedifferentiation and redifferentiation rates are modeled as non-linear functions of metabolic stress and lipid clearance:
\begin{equation}
k_{\mathrm{dediff}}(G, \mu_{\mathrm{fat}}) = k_{d0} \frac{G^2}{\theta_G^2 + G^2} (1 + \mu_{\mathrm{fat}}), \quad k_{\mathrm{rediff}}(\mu_{\mathrm{fat}}) = \frac{k_{r0}}{1 + 2\mu_{\mathrm{fat}}},
\end{equation}
with $k_{d0} = 0.03\text{ day}^{-1}$, $\theta_G = 140.0\text{ mg/dL}$, and $k_{r0} = 0.25\text{ day}^{-1}$.

\subsection{Autoimmune Surveillance Dynamics}

The autoimmune subsystem models clonal expansion of autoreactive effector T-cells, regulatory suppression by Tregs, and peripheral antigen presentation:
\begin{align}
\frac{dT_{\mathrm{eff}}}{dt} &= \left( \frac{\gamma_{\mathrm{eff}} \mathrm{Ag}}{\alpha_{\mathrm{Ag}} + \mathrm{Ag}} \right) T_{\mathrm{eff}} - \rho_{\mathrm{reg}} T_{\mathrm{reg}} T_{\mathrm{eff}} - d_{\mathrm{eff}} T_{\mathrm{eff}}, \label{eq:dT_eff} \\
\frac{dT_{\mathrm{reg}}}{dt} &= \gamma_{\mathrm{reg}} T_{\mathrm{reg}} - d_{\mathrm{reg}} T_{\mathrm{reg}} + u_{\mathrm{imm}}(t), \label{eq:dT_reg} \\
\frac{dC_{\mathrm{inflam}}}{dt} &= \eta_T T_{\mathrm{eff}} + \eta_G \max\left(0, G - 100.0\right) - d_{\mathrm{cyt}} C_{\mathrm{inflam}}, \label{eq:dC}
\end{align}
where autoantigen presentation $\mathrm{Ag} = 0.05 \beta_{\mathrm{func}} + 0.1 \max\left(0, G - 120.0\right)$ couples chronic metabolic hyperglycemia and ER stress to autoimmune activation.

\section{Type 2 Diabetes: Least Action Pulse Steering and Autonomous Remission}

\subsection{Bistability and the Separatrix Barrier}

When intra-organ fat accumulates ($\mu_{\mathrm{fat}} > 1.0$), the unforced metabolic subsystem ($\bm{u} \equiv 0$) undergoes a subcritical bifurcation, giving rise to two distinct stable attractors separated by a codimension-1 separatrix manifold $\mathcal{W}^s(\mathcal{X}_{\mathrm{saddle}})$:
\begin{itemize}
    \item \textbf{Healthy Homeostatic Attractor $\mathcal{A}_{\mathrm{norm}}$:} Characterized by normal fasting glycemia ($G \approx 95\text{--}100\text{ mg/dL}$), low insulin resistance, and robust functional beta-cell mass ($\beta_{\mathrm{func}} \approx 300\text{ mg}$, $\beta_{\mathrm{dediff}} \approx 0$).
    \item \textbf{Diabetic Pathological Attractor $\mathcal{A}_{\mathrm{diab}}$:} Characterized by severe fasting hyperglycemia ($G > 200\text{ mg/dL}$), elevated dedifferentiation ($\beta_{\mathrm{dediff}} \gg \beta_{\mathrm{func}}$), and blunted insulin secretion.
\end{itemize}

Conventional pharmacotherapy acts merely as a permanent external forcing term $u_{\mathrm{ext}}$ that artificially shifts the coordinates of $\mathcal{A}_{\mathrm{diab}}$ without altering the underlying attractor landscape. Consequently, termination of medication immediately returns the patient to severe hyperglycemia.

\subsection{Optimal Pulse Steering via Willems Dissipative Inequality}

\begin{definition}[Metabolic Storage Function]
In the framework of Willems dissipative dynamical systems \cite{Willems1972}, let $S(\bm{X}): \mathbb{R}^n \to \mathbb{R}^+$ be a non-negative, continuously differentiable storage function measuring geometric state departure from the target homeostatic attractor $\mathcal{X}_{\mathrm{norm}} = [100.0, \, 10.0, \, 300.0, \, 0.0]^T$:
\begin{equation}
S(\bm{X}) = \frac{1}{2} (\bm{X} - \mathcal{X}_{\mathrm{norm}})^T \bm{P} (\bm{X} - \mathcal{X}_{\mathrm{norm}}),
\end{equation}
where $\bm{P} = \mathrm{diag}\left( 50^{-2}, \, 10^{-2}, \, 100^{-2}, \, 100^{-2} \right)$ is a positive-definite metric weighting matrix.
\end{definition}

A therapeutic intervention profile $\bm{u}(t) = [u_{\mathrm{meta}}(t), \, D_{\mathrm{meal}}(t)]^T$ is dissipative with respect to supply rate $s(\bm{u}, \bm{y}) = \bm{y}^T \bm{u} - \gamma \|\bm{u}\|^2$ if the dissipation inequality holds:
\begin{equation}
\frac{dS(\bm{X})}{dt} \le s(\bm{u}, \bm{y}).
\end{equation}

To achieve permanent remission, the intervention profile must drive the trajectory from $\mathcal{A}_{\mathrm{diab}}$ across the saddle barrier $\mathcal{X}_{\mathrm{saddle}}$ into the basin of attraction $\mathcal{B}(\mathcal{A}_{\mathrm{norm}})$ within finite time $T^*$. The optimal steering protocol solves the Bolza optimal control problem:
\begin{equation}
\min_{\bm{u}(t), T^*} \mathcal{J} = \int_0^{T^*} \left[ \|\bm{X}(t) - \mathcal{X}_{\mathrm{norm}}\|_{\bm{Q}}^2 + \|\bm{u}(t)\|_{\bm{R}}^2 \right] dt + \Phi(\bm{X}(T^*)),
\end{equation}
subject to the state boundary condition $\bm{X}(T^*) \in \mathcal{B}(\mathcal{A}_{\mathrm{norm}})$.

Once the state enters $\mathcal{B}(\mathcal{A}_{\mathrm{norm}})$, external inputs can be completely extinguished ($\bm{u}(t) \equiv 0$ for $t > T^*$). The autonomous dissipative drift $\dot{S}(\bm{X}) \le -\alpha S(\bm{X})$ then guarantees exponential convergence to $\mathcal{A}_{\mathrm{norm}}$:
\begin{equation}
\lim_{t \to \infty} \bm{X}(t) = \mathcal{X}_{\mathrm{norm}}.
\end{equation}
This mathematical condition formalizes the biological reality of drug-free diabetes remission.

\subsection{Numerical Phase-Portrait Simulation Results}

We evaluated this optimal steering protocol across an 180-day simulation initialized from a patient in severe diabetic decompensation ($G(0) = 220.0\text{ mg/dL}$, $\beta_{\mathrm{func}}(0) = 120.0\text{ mg}$, $\beta_{\mathrm{dediff}}(0) = 180.0\text{ mg}$, $\mu_{\mathrm{fat}} = 1.5$):

\begin{table}[htbp]
\centering
\small
\caption{Trajectory comparison between optimal pulse intervention (remission) and insufficient intervention (relapse).}
\label{tab:t2d_results}
\begin{tabular}{lcccc}
\toprule
\textbf{Protocol Phase} & \textbf{Time (Days)} & \textbf{Intervention Status} & \textbf{Optimal Pulse (Remission)} & \textbf{Insufficient (Relapse)} \\
\midrule
Baseline Baseline & Day 0 & Unforced ($u=0$) & $220.00\text{ mg/dL}$ ($120.0\text{ mg}$) & $220.00\text{ mg/dL}$ ($120.0\text{ mg}$) \\
Active Window & Day 14--70 & Caloric/Incretin Pulse & Intensive ($\mu_{\mathrm{fat}} \to 0.05$) & 3-Week Sub-optimal \\
Intervention End & Day 70 & External Cessation & $73.63\text{ mg/dL}$ (Saddle Crossed) & $469.10\text{ mg/dL}$ (Sub-barrier) \\
Long-term Follow-up & Day 180 & Completely Unforced & $\mathbf{97.26\text{ mg/dL}}$ ($\mathbf{295.4\text{ mg}}$) & $\mathbf{589.01\text{ mg/dL}}$ ($\mathbf{0.4\text{ mg}}$) \\
\midrule
\textbf{Clinical Outcome} & - & - & \textbf{Permanent Remission} & \textbf{Irreversible Relapse} \\
\bottomrule
\end{tabular}
\end{table}

As summarized in Table \ref{tab:t2d_results}, the 8-week optimal pulse reduced intra-organ fat, triggered rapid redifferentiation of dormant beta-cells, and crossed the saddle barrier at Day 70 ($G = 73.63\text{ mg/dL}$). During the subsequent 110 days without medication or dietary restriction, the state autonomously stabilized at normal fasting glycemia ($97.26\text{ mg/dL}$), confirming that remission is an inescapable mathematical consequence of attractor basin steering. In contrast, the sub-optimal intervention failed to clear ectopic fat, leaving the state short of the separatrix, whereupon the trajectory relapsed into terminal glucotoxic failure ($589.01\text{ mg/dL}$, $\beta_{\mathrm{func}} = 0.40\text{ mg}$).

\section{Type 1 Diabetes: Pre-Symptomatic Screening and Geometric AID}

\subsection{Pre-Critical Sturm-Liouville Gap Softening}

In Type 1 diabetes, autoimmune beta-cell lysis is preceded by months of sub-clinical immune remodeling. We formalize this pre-symptomatic period using the local Koopman generator of the immune state $\bm{X}_{\mathrm{imm}} = [T_{\mathrm{eff}}, \, T_{\mathrm{reg}}, \, C_{\mathrm{inflam}}]^T$.

Over a moving window of duration $\Delta T$, we identify the continuous-time infinitesimal generator $\bm{L}_{\mathrm{imm}} \in \mathbb{R}^{3 \times 3}$:
\begin{equation}
\bm{L}_{\mathrm{imm}} = \frac{1}{\Delta t} \ln \left( \bm{Y}_1 \bm{Y}_0^T (\bm{Y}_0 \bm{Y}_0^T + \alpha \bm{I})^{-1} \right).
\end{equation}

\begin{definition}[Immune Restoring Gap]
The spectral Sturm-Liouville gap $\Omega_{1,\mathrm{immune}}$ of the autoimmune network is defined as:
\begin{equation}
\Omega_{1,\mathrm{immune}} \equiv -\max_{1 \le i \le 3} \mathrm{Re}\left( \mu_i(\bm{L}_{\mathrm{imm}}) \right).
\end{equation}
\end{definition}

In healthy self-tolerance, regulatory T-cells exert strong contractive damping on autoreactive clones, yielding a positive spectral gap ($\Omega_{1,\mathrm{immune}} \ge 0.10$). As the autoimmune network approaches the tipping point toward clonal runaway, critical slowing down induces spectral gap softening:
\begin{equation}
\Omega_{1,\mathrm{immune}} \to 0^+.
\end{equation}

In our numerical simulation spanning 365 days of progressive insulitis:
\begin{itemize}
    \item The spectral gap $\Omega_{1,\mathrm{immune}}$ softened below the critical warning threshold ($0.02$) at **Day 84.9**.
    \item Macroscopic clinical diabetes ($G > 180.0\text{ mg/dL}$) manifested only at **Day 234.3**, when functional beta-cell mass had declined below $15\%$.
    \item The operator-theoretic framework achieved an autonomous pre-symptomatic early warning lead time of:
    \begin{equation}
    \Delta t_{\mathrm{lead}} = 234.3 - 84.9 = \mathbf{149.4\text{ days}} \quad (\approx \mathbf{5\text{ months}}).
    \end{equation}
\end{itemize}
This multi-month lead time provides a decisive therapeutic window for targeted tolerance-reinduction therapies (such as low-dose IL-2 or antigen-specific nanoparticle pulses) prior to irreversible islet devastation.

\subsection{Raychaudhuri-Informed Artificial Pancreas (AID) Control}

For patients with established beta-cell loss, external closed-loop automated insulin delivery (AID) is mandatory. Conventional Proportional-Integral-Derivative (PID) and linear Model Predictive Control (MPC) algorithms respond purely to scalar glucose deviations, frequently causing postprandial hyperglycemic overshoots followed by delayed hypoglycemic undershoots.

To overcome phase lag, we formulate a geometry-informed controller based on the discrete Raychaudhuri convergence scalar $\theta_k$:
\begin{equation}
\theta_k \equiv \mathrm{Tr}\left( \bm{L}_{\mathrm{meta}, k} \right),
\end{equation}
derived from delay-coordinate embeddings $\bm{Z}_k = [G_k, \, G_{k-\tau}, \, \dots, \, G_{k-(d-1)\tau}]^T$ ($d = 4$, $\tau = 3\text{ min}$).
\begin{itemize}
    \item \textbf{Satisfied (Contractive Dynamics):} $\theta_k < 0$. The phase-space bundle contracts, indicating stable closed-loop regulation.
    \item \textbf{Unsatisfied (Divergent Expansion):} $\theta_k > 0.04$. The bundle volume expands rapidly, signaling incipient glycemic runaway due to unannounced meals or exercise.
\end{itemize}

The Raychaudhuri-AID controller augments standard basal-bolus feedback with an anticipatory geometric pulse:
\begin{equation}
u_{\mathrm{AID}}(k) = u_{\mathrm{PID}}(k) + \mathcal{H}(\theta_k - 0.04) \cdot \Phi_{\mathrm{anticip}}(\theta_k, G_k),
\end{equation}
where $\mathcal{H}$ is the Heaviside step and $\Phi_{\mathrm{anticip}}$ delivers an immediate preemptive micro-bolus before peripheral glucose sensors register large excursions.

In an acute in-silico challenge test (ingestion of $75\text{ g}$ carbohydrate bolus):
\begin{itemize}
    \item Standard PID control permitted postprandial glucose to peak at $175.64\text{ mg/dL}$, followed by a delayed undershoot to $95.32\text{ mg/dL}$.
    \item The Raychaudhuri-AID algorithm blunted the peak excursion to $168.93\text{ mg/dL}$ and completely stabilized postprandial recovery at a safe nadir of $\mathbf{105.23\text{ mg/dL}}$, entirely eliminating reactive hypoglycemia.
\end{itemize}

\section{Structural Bifurcation Triage and Point of No Return}

A critical requirement in clinical deployment is the rigorous triage between patients eligible for curative attractor steering and those requiring permanent technological replacement.

Total beta-cell mass is partitioned into functional, dedifferentiated, and irreversibly destroyed compartments:
\begin{equation}
\beta_{\mathrm{total}} = \beta_{\mathrm{func}}(t) + \beta_{\mathrm{dediff}}(t) + \beta_{\mathrm{dead}}(t).
\end{equation}
Let $\rho_{\mathrm{dead}} = \beta_{\mathrm{dead}} / \beta_{\mathrm{total}} \in [0, 1]$ denote the fraction of permanently lost islet mass.

\begin{theorem}[Fold Bifurcation and Remission Threshold]
In the unforced autonomous metabolic subsystem ($\bm{u} \equiv 0$), there exists a critical irreversible loss threshold $\rho_{\mathrm{crit}} \in (0, 1)$ such that:
\begin{enumerate}
    \item For $\rho_{\mathrm{dead}} < \rho_{\mathrm{crit}}$, the system possesses a non-empty healthy homeostatic attractor $\mathcal{A}_{\mathrm{norm}}$ and a distinct saddle $\mathcal{X}_{\mathrm{saddle}}$. The system is \textbf{topologically reversible}.
    \item At $\rho_{\mathrm{dead}} = \rho_{\mathrm{crit}}$, the healthy attractor $\mathcal{A}_{\mathrm{norm}}$ and the saddle $\mathcal{X}_{\mathrm{saddle}}$ coalesce in a codimension-1 saddle-node (fold) bifurcation.
    \item For $\rho_{\mathrm{dead}} > \rho_{\mathrm{crit}}$, $\mathcal{A}_{\mathrm{norm}}$ ceases to exist in real phase space. The diabetic state $\mathcal{A}_{\mathrm{diab}}$ becomes a global attractor. The system is \textbf{topologically irreversible}.
\end{enumerate}
\end{theorem}

\begin{proof}
Consider the quasi-steady-state nullclines of Equations (\ref{eq:dG})--(\ref{eq:dbeta_f}) with respect to functional mass $\beta_{\mathrm{func}}$. Solving $dG/dt = 0$ yields $G = R_0 / (E_{G0} + S_I I)$. Substituting $I = (\sigma_{\mathrm{ins}} \beta_{\mathrm{func}} G^2) / [k_e(\alpha_G + G^2)]$ defines a monotonic algebraic relation $G = \mathcal{G}(\beta_{\mathrm{func}})$. 

Substituting $\mathcal{G}(\beta_{\mathrm{func}})$ into the steady-state beta-cell turnover condition $F(\beta_{\mathrm{func}}; \rho_{\mathrm{dead}}) \equiv f_{\mathrm{turn}}(\mathcal{G}(\beta_{\mathrm{func}})) \beta_{\mathrm{func}} + k_{\mathrm{rediff}} \beta_{\mathrm{dediff}} - k_{\mathrm{dediff}} \beta_{\mathrm{func}} = 0$ yields a cubic characteristic equation in $\beta_{\mathrm{func}}$. When $\beta_{\mathrm{dead}} \le \rho_{\mathrm{crit}} \beta_{\mathrm{total}}$, the discriminant $\Delta(\rho_{\mathrm{dead}})$ is strictly positive, yielding three real roots: $\beta_{\mathrm{norm}} > \beta_{\mathrm{saddle}} > \beta_{\mathrm{diab}}$. 

As $\rho_{\mathrm{dead}}$ increases, the available pool $\beta_{\mathrm{func}} + \beta_{\mathrm{dediff}} = (1 - \rho_{\mathrm{dead}})\beta_{\mathrm{total}}$ contracts, shifting the upper nullcline downward until the discriminant vanishes:
\begin{equation}
\Delta(\rho_{\mathrm{crit}}) = 0.
\end{equation}
At this point, $\partial F / \partial \beta_{\mathrm{func}} = 0$ and $\partial^2 F / \partial \beta_{\mathrm{func}}^2 \neq 0$, satisfying Sotomayor's conditions for a fold bifurcation. For $\rho_{\mathrm{dead}} > \rho_{\mathrm{crit}}$, $\Delta < 0$, eliminating real solutions for $\mathcal{A}_{\mathrm{norm}}$.
\end{proof}

\begin{remark}[Clinical Triage Protocol]
Theorem 5.1 provides the definitive mathematical rule for clinical triage:
\begin{itemize}
    \item \textbf{Category I ($\rho_{\mathrm{dead}} < \rho_{\mathrm{crit}}$, Reversible Phase):} Dedifferentiation dominates ($\beta_{\mathrm{dediff}} \gg \beta_{\mathrm{dead}}$). The patient is assigned to the \textbf{Finite Least Action Remission Protocol} (Section 3).
    \item \textbf{Category II ($\rho_{\mathrm{dead}} \ge \rho_{\mathrm{crit}}$, Point of No Return):} Apoptosis/destruction dominates. Curative attractor steering is mathematically impossible; the patient is immediately triaged to the \textbf{Raychaudhuri-AID Autonomous Artificial Pancreas Protocol} (Section 4.2).
\end{itemize}
\end{remark}

\section{Physical Isomorphisms with Holographic Self-Regulation}

The mathematical structures governing biological metabolic-immune homeostasis exhibit rigorous structural isomorphisms with the holographic cosmological self-regulation framework established in Paper VI \cite{Yoshikai2026f}:

\begin{table}[htbp]
\centering
\small
\caption{Scale-invariant isomorphisms between cosmological self-regulation and biophysical homeostasis.}
\label{tab:isomorphisms}
\begin{tabularx}{\textwidth}{>{\raggedright\arraybackslash}p{4.2cm} >{\raggedright\arraybackslash}p{5.5cm} >{\raggedright\arraybackslash}X}
\toprule
\textbf{Cosmological / Holographic Structure (Paper VI)} & \textbf{Metabolic-Immune Network (Paper VIII)} & \textbf{Shared Mathematical Universality Class} \\
\midrule
Vacuum stability against decay ($\delta B_{\mathrm{CDL}} \ge 0$) & Self-tolerance preservation against autoimmunity ($\Omega_{1,\mathrm{immune}} > 0$) & Positive spectral gap in self-adjoint dissipative generator \\
Moduli gradient flow on Weil-Petersson space & Metabolic trajectory convergence to homeostatic attractor $\mathcal{A}_{\mathrm{norm}}$ & Asymptotic stability via positive-definite storage functions $\dot{S} \le 0$ \\
$T\bar{T}$ Dirichlet cutoff boundary & Beta-cell membrane refractory gating and glycemic thresholds & Finite-boundary truncation in nonlinear boundary-value problems \\
Hawking radiation and bulk open-closed mode transduction & Irreversible hepatic/renal metabolic waste clearance & Entropy export to ambient reservoir in open Hamiltonian systems \\
Event horizon formation and singularity encapsulation & Irreversible beta-cell apoptosis ($\beta_{\mathrm{dead}}$) and scar formation & Non-invertible topological information drop from state space \\
\bottomrule
\end{tabularx}
\end{table}

These structural parallels demonstrate that living matter and gravitational spacetime are governed by the same underlying hydrodynamic central limit theorem (Theorem 4.2 in \cite{Yoshikai2026f}): coarse-grained macroscopic stability in open systems unconditionally contracts onto identical nonlinear dissipative universality classes.

\section{Conclusion}

We have established a unified dissipative systems framework for the prevention, remission, and autonomous management of diabetes. By reframing metabolic and autoimmune pathologies as non-equilibrium attractors on a multi-scale manifold:
\begin{enumerate}
    \item We demonstrated that Type 2 diabetes remission is achievable through finite-duration, least-action pulse interventions that steer the state across the separatrix into the homeostatic basin of attraction, enabling drug-free glycemic stability ($97.26\text{ mg/dL}$).
    \item We formulated an autonomous early warning screening method for Type 1 diabetes via Sturm-Liouville gap softening ($\Omega_{1,\mathrm{immune}} \to 0$), detecting pre-symptomatic phase transitions $149.4\text{ days}$ prior to macroscopic clinical hyperglycemia.
    \item We engineered a geometry-informed artificial pancreas algorithm utilizing Raychaudhuri bundle expansion scalars $\theta_k$, preempting postprandial spikes and maintaining postprandial glucose strictly above hypoglycemic thresholds ($105.23\text{ mg/dL}$).
    \item We proved a fold bifurcation theorem that provides an analytical triage between reversible dedifferentiation suitable for curative steering and irreversible destruction requiring robotic augmentation.
\end{enumerate}

These findings transform diabetes care from perpetual palliative stabilization into proactive, geometry-guided curative medicine.

\section*{Acknowledgments}
The author acknowledges the use of the artificial intelligence model Gemini (developed by Google) for assistance with multi-scale endocrine-immune dynamical modeling, numerical simulation code development, peer-review simulation, and manuscript preparation. All core theoretical concepts, formal mathematical formulations, and physiological conclusions remain the sole responsibility of the author.

\section*{Data and Code Availability}

The complete Python simulation scripts, differential equation solvers, and numerical trajectory datasets supporting this study are available in the open-source repository under the MIT License at: \url{https://github.com/sukuyoshi/metabolic_immune_dynamics}.

\begin{thebibliography}{99}

\bibitem{DeFronzo2009}
R.~A.~DeFronzo,
\newblock From the triumvirate to the ominous octet: A new paradigm for the treatment of type 2 diabetes mellitus,
\newblock \emph{Diabetes}, 58(4):773--795, 2009.

\bibitem{Kahn2006}
S.~E.~Kahn, R.~L.~Hull, and K.~M.~Utzschneider,
\newblock Mechanisms linking obesity to insulin resistance and type 2 diabetes,
\newblock \emph{Nature}, 444(7121):840--846, 2006.

\bibitem{Eisenbarth1986}
G.~S.~Eisenbarth,
\newblock Type I diabetes mellitus: A chronic autoimmune disease,
\newblock \emph{New England Journal of Medicine}, 314(21):1360--1368, 1986.

\bibitem{Atkinson2014}
M.~A.~Atkinson, G.~S.~Eisenbarth, and A.~W.~Michels,
\newblock Type 1 diabetes,
\newblock \emph{The Lancet}, 383(9911):69--82, 2014.

\bibitem{Talchai2012}
C.~Talchai, S.~Xuan, H.~V.~Lin, L.~Sussel, and D.~Accili,
\newblock Pancreatic $\beta$ cell dedifferentiation as a mechanism of diabetic $\beta$ cell failure,
\newblock \emph{Cell}, 150(6):1223--1234, 2012.

\bibitem{Accili2019}
D.~Accili,
\newblock The direct and indirect effects of insulin on pancreatic $\beta$ cells,
\newblock \emph{Cell Metabolism}, 29(4):795--797, 2019.

\bibitem{Lean2018}
M.~E.~Lean, W.~S.~Leslie, A.~C.~Barnes, N.~Brosnahan, G.~Khomami, C.~McCombie, L.~McIntyre, R.~McSorley, R.~Morgan, R.~Taylor, et~al.,
\newblock Primary care-led weight management for remission of type 2 diabetes (DiRECT): An open-label, cluster-randomised trial,
\newblock \emph{The Lancet}, 391(10120):541--551, 2018.

\bibitem{Taylor2018}
R.~Taylor, A.~Al-Mrabeh, S.~Zhyzhneuskaya, C.~Peters, A.~C.~Barnes, B.~S.~Aribisala, K.~G.~Hollingsworth, J.~Mathers, N.~Sattar, and M.~E.~Lean,
\newblock Remission of human type 2 diabetes, requires decrease in liver and pancreas fat content but is dependent upon capacity for $\beta$ cell recovery,
\newblock \emph{Cell Metabolism}, 28(4):547--556, 2018.

\bibitem{Taylor2019}
R.~Taylor,
\newblock Calorie restriction for long-term remission of type 2 diabetes,
\newblock \emph{Clinical Medicine}, 19(1):37--42, 2019.

\bibitem{Ziegler2013}
A.-G.~Ziegler, M.~Rewers, O.~Simell, T.~Simell, J.~Lempainen, {\AA}.~Steck, C.~Winkler, J.~Ilonen, D.~Dunger, J.~M.~She, et~al.,
\newblock Seroconversion to multiple islet autoantibodies and risk of progression to type 1 diabetes in children,
\newblock \emph{JAMA}, 309(23):2473--2479, 2013.

\bibitem{Insel2015}
R.~A.~Insel, J.~L.~Dunne, M.~A.~Atkinson, J.~L.~Chiang, D.~Dabelea, P.~A.~Gottlieb, C.~J.~Greenbaum, K.~C.~Herold, J.~P.~Krischer, {\AA}.~Lernmark, et~al.,
\newblock Staging presymptomatic type 1 diabetes: A scientific statement of JDRF, the Endocrine Society, and the American Diabetes Association,
\newblock \emph{Diabetes Care}, 38(10):1964--1974, 2015.

\bibitem{Willems1972}
J.~C.~Willems,
\newblock Dissipative dynamical systems part I: General theory; Part II: Linear systems with quadratic supply rates,
\newblock \emph{Archive for Rational Mechanics and Analysis}, 45(5):321--393, 1972.

\bibitem{Prigogine1977}
I.~Prigogine,
\newblock \emph{Self-Organization in Nonequilibrium Systems: From Dissipative Structures to Order through Fluctuations},
\newblock John Wiley \& Sons, New York, 1977.

\bibitem{Scheffer2009}
M.~Scheffer, J.~Bascompte, W.~A.~Brock, V.~Brovkin, S.~R.~Carpenter, V.~Dakos, H.~Held, E.~H.~van~Nes, M.~Rietkerk, and G.~Sugihara,
\newblock Early-warning signals for critical transitions,
\newblock \emph{Nature}, 461(7260):53--59, 2009.

\bibitem{Topp2000}
B.~Topp, K.~Promislow, G.~deVries, R.~M.~Miura, and D.~T.~Finegood,
\newblock A model of $\beta$-cell mass, insulin, and glucose kinetics: Pathways to diabetes,
\newblock \emph{Journal of Theoretical Biology}, 206(4):605--619, 2000.

\bibitem{Raychaudhuri1955}
A.~Raychaudhuri,
\newblock Relativistic cosmology. I,
\newblock \emph{Physical Review}, 98(4):1123--1126, 1955.

\bibitem{Yoshikai2026f}
R.~Yoshikai,
\newblock The autopoietic universe: Non-equilibrium self-regulation, holographic dissipative systems, and moduli gradient flow,
\newblock Zenodo, 2026. doi:\href{https://doi.org/10.5281/zenodo.24156018}{10.5281/zenodo.24156018}.

\bibitem{Yoshikai2026g}
R.~Yoshikai,
\newblock Operator-theoretic early warning diagnostics for non-equilibrium phase transitions: Predicting lethal cardiac cascades via Sturm-Liouville gap softening and Raychaudhuri expansion,
\newblock Zenodo, 2026. doi:\href{https://doi.org/10.5281/zenodo.24156142}{10.5281/zenodo.24156142}.

\end{thebibliography}

\end{document}
